% Folder 83, batch 04 (archivists' pages 61-80).
% Pass: Opus 5.5 (claude-opus-5-5), 2026-10-10 — first pass, unchecked against the pages by a human.
\documentclass[11pt,a4paper]{article}
\input{../preamble/grothendieck}

\folder{83}
\batch{4}
\pages{61}{80}
\foldertitle{[Icosaèdres, bi-icosaèdres et algèbres d'Azumaya de rang 4] : notes manuscrites (1982-1983, 1986, s.d.).}
\dating{1982-1986}
\watermark{Édition de démonstration}

\begin{document}

\page{61}
\note{la page s'ouvre sur une nouvelle mise en place, sans renvoi explicite aux pages précédentes ; elle est très raturée.}
Soit $M$ alg. d'Azumaya sur corps $\mathbb{F}_p$, $p$ premier impair, et $C^*$ \uncertain{l'ens.} des éléments nilpotents non nuls de $M$.
Alors \struck{\ill{}} $C^*$ se \uncertain{divise} \uncertain{naturellement} en deux $C^*_i$, \struck{\ill{}}
\[
  \lambda\, C^*_i = C^*_i \quad\text{ssi}\quad \lambda \in k^{*2} \qquad (\text{pour } \lambda\in k^*) .
\]
\note{sous $k^{*2}$ est inscrit, en petit, « $\simeq k^*/\pm 1$ » ; lecture douteuse.}

Sur $C^*$, il y a une str. \uncertain{canon.} de \uncertain{bi-cônes}, \struck{il y a une \ill{} \ill{} sur chaque $C_i$, \ill{} des \ill{} \ill{} des cônes \uncertain{orientés} \ill{} \ill{} \uncertain{multi-cônes}}, et les opérations $\overline{T}_\lambda$ ($\lambda\in k^*$) \uncertain{induisent} les \uncertain{homothéties} \ill{} $\lambda^2$. Le groupe $Gp_M(\struck{\ill{}}) \simeq M^*/k^*$ \struck{\ill{} \ill{}} \struck{\ill{} $C_i$ \ill{} \ill{}} \struck{groupe des \uncertain{cônes}} \add{des cônes} bi-cônes \uncertain{orientés}, il a un sous-groupe d'indice $2$ : les \ill{} qui \uncertain{préservent} les \ill{} \ill{} \uncertain{mieux}. \struck{Le groupe}

\struck{$Gp_M\times k^*$ est le groupe des \uncertain{automorphismes}} \struck{des} \struck{\ill{}} \struck{qui \uncertain{préserve} \ill{}} des \uncertain{structures} de \ill{} $C_i$ \struck{et l'un $E$} : Considérons \ill{} \ill{} \ill{} \ill{} $C_i$, \ill{} \ill{} \ill{} les \uncertain{bi-cônes} \ill{} par les \uncertain{transformations} $\overline{T}_{\lambda^2}(c_i)$ \add{(\ill{} \ill{} les $\overline{T}_{\lambda^2}(c_i)$)}.

$\overline{T}_\lambda$ \emph{\uncertain{transporte}} \emph{des structures} : \struck{\ill{}} \add{\ill{}}
$\{a,b\}$ est une \ill{} préservée par $T_\lambda(c_i)$ ssi \ill{} $\{T_\lambda a, T_\lambda b\}$ \uncertain{dans} est une \ill{} par $c_i$ \emph{\uncertain{et non pas}} $\{a, T_\lambda b\}$ \ill{} \ill{} \struck{\ill{} \ill{} \ill{}} \add{si $\lambda^2 = -1$, on trouve} donc $T_\lambda(c_i) = c_i$, mais $T'_\lambda$ est une \uncertain{autre}…

\marginal{il y a une \ill{} \uncertain{classique} (cf. p.~27\,?), \ill{} \ill{} \ill{} \ill{} \ill{} bien \uncertain{mystérieux}\,! \ill{} \ill{} \ill{} \ill{}}
\note{note marginale écrite en oblique dans l'angle gauche ; presque entièrement illisible. Le renvoi « p.~27 » est de l'auteur ; lecture douteuse.}

\page{62}
de ce qui \emph{renverse} l'orientation. \struck{Ce}
Donc la structure $E$, dans le cas de l'iso. $p=5$, où $k^{*2} = \{\pm 1\}$, est la \add{icosaédrale} (une orientation de bi-icosaèdre). Ceci dit, \ill{} \ill{} $Gp\times k^*$ opérant \ill{} le groupe des \ill{} \ill{} \ill{} \uncertain{structure} de multi-cônes.
\note{dans « $k^{*2}$ » l'exposant est surchargé de traits ; lecture douteuse.}

On trouve une représentation \uncertain{exceptionnelle}, remarquable, des \uncertain{autres} \ill{} $Gp$. \struck{par les} Considérons les \uncertain{couples} universels (« définis sur $\mathbb{Z}$ ») qui sont ($\rho_f$ d'ordre $3$, \add{$\det\rho_f = 1$, $\operatorname{Trace}\rho_f = -1$}) et $\rho_s$ \struck{\ill{}} \emph{unipotent} ($\det\rho_s = 1$, $\operatorname{Tr}\rho_s = 2$), lesquelles $\rho_s$ soit unipotent, dans Alg d'Azumaya engendrée par
\[
  \rho_a = \sigma \text{ d'ordre } 4, \qquad \rho_f \text{ d'ordre } 3,
\]
\[
  \left\{
  \begin{array}{ll}
    \det\sigma = 1 , & \det\rho_f = 1 \\
    \operatorname{Tr}\sigma = 0 , & \operatorname{Tr}\rho_f = -1
  \end{array}
  \right.
\]
avec
\[
  \sigma . \rho_f = \rho_s^{-1} \qquad \rho_s \text{ unipotent i.e. } \det\rho_s = 1,\ \operatorname{Tr}\rho_s = 2 .
\]
\note{dans « $\sigma.\rho_f = \rho_s^{-1}$ » un symbole entre $=$ et $\rho_s$ est noirci et biffé.}

\marginal{NB l'homothétie $-1$ \uncertain{inverse} les deux \ill{} d'\uncertain{isotropie}, la structure ; il \ill{} deux $C_i$ ssi $-1 \in k^{*2}$ — \ill{} \ill{} un autre \ill{} \uncertain{orientés} de structure \ill{}. Je \ill{} sur l'expression $C^*$ des \ill{} \ill{} \ill{} (dans les \ill{} \ill{}) \ill{} \ill{}}
\note{note marginale écrite en oblique sur la partie gauche de la page, en travers du texte ; lecture très incomplète.}

\page{63}
On trouve une représentation \uncertain{fidèle} de $SL(2,\mathbb{Z})/\pm 1$ dans le groupe \ill{} par $\sigma$ et $\rho_f$ mod \uncertain{le centre} de l'alg. d'Azumaya.) \struck{\ill{} \ill{} \ill{}} Considérons \struck{\ill{}} la représentation de la \ill{} \ill{} \uncertain{bi-}\ill{}-universelle, dans les \add{\ill{}} \struck{\ill{}} \ill{} ($\operatorname{Tr} u = 1$, \ill{}) engendrée par les \struck{\ill{}} $\rho_s$. \struck{En \ill{}}

En réduisant mod $n$, je dis que ça se factorise par $SL(2,\mathbb{Z}/n)/\pm 1$ (tautologique, tel que j'ai présenté \ill{} choses), et la représentation \ill{} est celle que j'ai décrite tantôt (mais qui dit \ill{} un sens, en supposant $n$ premier…).
\note{dans « $SL(2,\mathbb{Z}/n)$ » le $n$ est écrit par-dessus un autre symbole biffé.}

\marginal{Les \uncertain{relations} sont \uncertain{les mêmes} \struck{\ill{}} \ill{} \ill{} \ill{}. Les \uncertain{générateurs} \ill{} : paires d'\uncertain{unipotents} $(\rho_s, \sigma\rho_s\sigma^{-1})$, $\rho_{s'}$ où $s' = \sigma(s)$, i.e. paires d'unipotents $(u,u')$ tels que $\operatorname{Tr} uu' = -1$ \ill{} $-\sigma\rho_s\sigma$ \ill{}}
\note{note marginale en oblique sur le bord gauche, en partie biffée.}

\page{64}
Soit $A$ une algèbre d'Azumaya de rang $4$ sur un anneau $k$, et supposons que $M$ \struck{$M$} soit « triviale », plus précisément, qu'elle provienne d'un $k$-module projectif de rang $2$ $E$, tel que le module inversible $\det E$ soit trivial (i.e. le \struck{groupe \ill{} \ill{}} $Gp(1)_k$-torseur \uncertain{déduit} \uncertain{venant} de $M$ que \ill{} \uncertain{remonte} en un $Sl(2)_k$-torseur.) Mais pour le moment, supposons seulement que $M$ soit triviale i.e.
\[
  M \simeq \operatorname{End}(E) .
\]
\note{« Soit $A$ » au début, puis $M$ dans toute la suite.}

Soit $u\in M$ un élément \struck{\ill{}} nilpotent dégénéré de carré nul. Alors $\operatorname{Im} u = \operatorname{Ker} u$ définit une « droite » $L\subset M$, et $D = ku \simeq L\otimes L'$, où $L'$ est l'orthogonal de $L$ dans $\check{E} \simeq E\otimes \struck{\ill{}}^{-1}$ (où $\underline{\omega} = \omega(E)$).
\struck{Dans} En fait, \ill{} \ill{} l'identification précédente (définie au \uncertain{signe} près, \uncertain{conventions de signes}) \ill{}
\[
  L' = L\otimes \underline{\omega}^{-1} \qquad\text{donc}
\]
\[
  D \simeq L^{\otimes 2}\otimes\underline{\omega}^{-1}
  \quad\text{i.e.}\quad
  \underline{\omega} \xrightarrow[\sim]{\varphi_L} \underline{\operatorname{Hom}}(L^{\otimes 2}, D_L)
\]
\note{sous le premier $D$ est inscrit l'indice $L$ : $D_L$.}
[Ceci, au moins \uncertain{pour} \uncertain{fait} la conv. de signe].
Le choix d'une base $e\in\underline{\omega}$, d'où $e^{-1}\in\underline{\omega}^{-1}$, définit un iso.
\[
  \varphi_L^e : \struck{\ill{}} \; L^{\otimes 2} \xrightarrow{\ \sim\ } D_L
\]
Si on remplace $e$ par $\lambda e$, $e^{-1}$ est remplacé par $\lambda^{-1}e^{-1}$ et on trouve
\[
  \varphi_L^{\lambda e} = \lambda\, \varphi_L^{e} .
\]
Quand $e$ est choisi (\uncertain{str.} symplectique sur $E$)

\page{65}
d'où l'iso $\varphi^e$, et si on suppose $L$ trivial i.e. $L^*$ non vide, \struck{\ill{}} \ill{} \ill{} pour $D$, correspondant \ill{} \struck{des} \uncertain{l'hyp.} qui \ill{} \ill{} \ill{} \ill{} $ku$) dans $D^*$ il y a une classe d'éléments modulo $k^{*2}$, privilégiée : celle formée par les $\varphi^e(x\otimes x)$, pour $x\in L^*$. Appelons-la $D^*(e)$ : \struck{On aura $D^*(\lambda e) = \lambda D^*(e)$.}
\[
  D^*(e) = \{\, \varphi_L^e(x\otimes x) \mid x\in L^* \,\} \qquad\text{pour } D = D_L
\]
donc
\[
  D^*(\lambda e) = \lambda D^*(e)
\]
et bien sûr $\lambda^2 D^*(e) = D^*(e)$, puisque $D^*(e) = k^{*2}u$ pour $u\in D^*(e)$.

\struck{La subdivision de $D^*$ \ill{} \ill{}}
On peut dire que \add{\uncertain{un iso.} \uncertain{can.}} \uncertain{quant} le $k^*$-torseur $D^*$, \ill{} \ill{}
\[
  D^*/k^{*2} \simeq \underline{\omega}^*/k^{*2} .
\]

\struck{\ill{}} Supposons que ${}_2\mathrm{Pic}(k) = 0$, donc \uncertain{pour} toute \struck{droite $L$ de} $u$ dans le cône $C^*$ des éléments de $M$ nilpotents dégénérés de carré nul, (i.e. tels que $\operatorname{Tr} u = \det u = 0$) la droite $L_u$ \uncertain{est} non vide. Alors

On définit $\struck{\ill{}}$
\[
  C^*(e) = \bigcup_{D} D^*(e)
\]
et on a ceci :
\[
  \left\{
  \begin{array}{l}
    C^*(\lambda e) = \lambda C^*(e) \\
    \lambda^2 C^*(e) = C^*(e)
  \end{array}
  \right.
\]
donc $C^*(e)$ ne dépend que de $e$ mod $k^{*2}$ \struck{\ill{}}.

\page{66}
et on trouve que
\[
  C^* = \coprod_{e\in\underline{\omega}^*/k^{*2}} C^*(e)
\]
\note{sous l'indice $\underline{\omega}^*/k^{*2}$, souligné d'un trait ondulé : « torseur sous $k^*/k^{*2}$ ».}
partition qui est permutée transitivement par les homothéties \add{de $M$}, \struck{L'\ill{}} que $\lambda C^*(e) = C^*(\lambda e)$. L'ens. des classes de cette partition est \struck{\ill{} \ill{}} de façon naturelle \uncertain{(via} la relation d'homothéties) un torseur sous $k^*/k^{*2}$, et ce torseur est can. iso à $\underline{\omega}^*/k^{*2}$.

Cette subdivision de $C^*$ en classes ne dépend que de $M$, et pas du \uncertain{choix} de $E$ (avec $\det E\simeq k$) et d'un iso $\operatorname{End}(E)\simeq M$. On \uncertain{va} \ill{} plus loin, \struck{\ill{} \ill{} On}.

On \struck{\ill{}} a une application surjective
\[
  \varphi^e : E^* \longrightarrow C^*(e), \qquad x\longmapsto \varphi^e(x\otimes x)
\]
satisfaisant
\[
  \left\{
  \begin{array}{l}
    \varphi^e(\lambda x) = \lambda^2 \varphi^e(x) \\
    \varphi^{\lambda e}(x) = \lambda\, \varphi^e(x)
  \end{array}
  \right.
\]

\marginal{On \ill{} ${}_2\mathrm{Pic}(k) = 0$ \ill{} \ill{} (\ill{} \ill{}) \ill{} \ill{} \ill{} \ill{} d'une \ill{} \ill{} \ill{}}
\note{note marginale oblique, reliée par une flèche au texte ; lecture très incomplète. Un signe de vérification (double trait) est tracé dans la marge gauche.}

\page{67}
Ainsi, \struck{\ill{}} \add{chaque} classe $C^*_i$ \add{($= C^*(e)$)} \struck{($e$ fixé)} est munie \struck{\ill{}} d'une structure \add{(indépendante du choix de $e$)} \uncertain{tri-angulaire} \struck{\ill{}} \add{orientée} canonique, et en plus d'une \ill{}-orientation \add{(dépendant a priori du choix de $e$)} qu'il faudrait expliciter. Le groupe projectif de $M$, $M^*/k^*$ ($\simeq \operatorname{End}(E)/k^*$) permute \struck{deux} entre elles les $C^*_i$, \add{(par transport des structures)} en respectant les structures \uncertain{et} orientations.

\struck{On \ill{}} \struck{Les \uncertain{autres}} On fera attention \struck{qui \ill{}} que pour $u\in GL(E)$, définissant $\bar u\in \struck{GL(E)} \, GL(E)/k^*\simeq M^*/k^*$ \struck{\ill{}} \add{(opérant \uncertain{dans} $C^*$)}, l'automorphisme de $M^*$ \struck{\ill{}} qu'il définit, est donné par
\[
  \Theta \longmapsto u(\Theta) = u\,\Theta\, u^{-1} .
\]
Ceci est compatible avec l'application
\[
  \varphi^e : E^* \xrightarrow{\ \sim\ } C^*(e) \subset C^*
\]
(avec $GL(E)$ opérant sur $E^*$ et sur $C^*$) que sur le sous-groupe $SL(E)$. De façon précise, on aura
\[
  \boxed{\ \Theta(\varphi^e(x)) = \det\Theta\,\bigl(\varphi^e(\Theta(x))\bigr)\ }
\]
\[
  \Bigl(\, \Theta(x\overset{e}{\otimes} x) = \Theta(x)\overset{\Theta(e)}{\otimes}\Theta(x)
  = \Theta(x)\overset{\det\Theta.e}{\otimes}\Theta(x)
  = \det\Theta\,\bigl(\Theta(x)\overset{e}{\otimes}\Theta(x)\bigr) \Bigr)
\]
\struck{Dans le groupe}

\marginal{\emph{Attention}}
\note{« Attention » souligné, dans la marge gauche, devant un double trait vertical qui encadre l'alinéa « On fera attention… ».}

\marginal{\ill{} \ill{} \uncertain{ici} \ill{} il faut \ill{} $C^*$ \ill{} \ill{} \ill{} $\operatorname{Tr} u = 0$ \ill{} \ill{} \ill{} \ill{}}
\note{note marginale oblique dans l'angle gauche ; lecture très incomplète.}

\page{68}
\emph{Remarques} Indépendamment de l'hyp. sur ${}_2\mathrm{Pic}(k)$ et sur $M$, on peut regarder $C^*$, et la structure de « graphe » sur $C^*$ définie par la relation $\operatorname{Tr}(uv) = -1$, d'où structure triangulaire. Mais

1°) $C^*$ n'est non vide que si $M$ est triviale, plus précisément si $M\simeq\operatorname{End}(E)$ avec $E = k\oplus L$ ($L$ module inv.)

\struck{On a \uncertain{dit}}
2°) \struck{Il \ill{} Je} L'un des côtés est non vide ssi $M \simeq M_2(k)$. En effet, si $\exists$ côté non vide $\{u,v\}$ \struck{\uncertain{Trace}} avec $u,v\in C^*$, $\operatorname{Tr} uv = -1$, les \uncertain{relations} de conjugaison nous fournissent un \uncertain{homomorphisme} $M_2 \to M$, transformant $(u_0,v_0)$ \uncertain{couple} type en $(u,v)$ (\uncertain{donc} \uncertain{on a} \ill{} conditions\,!). Bon aussi que $u_0$ en $u$, $v_0$ en $v$ dans $M$, satisfaisant les \ill{}

Donc on a \uncertain{indirectement} si on beaucoup \ill{} $M\simeq\operatorname{End}(E)$, un supplémentaire $E$ \ill{} \ill{} \ill{} iso à deux \ill{} \emph{\uncertain{libre}}.
\note{la fin de la page est d'une lecture très incertaine ; le fil reprend p.~69.}

\page{69}
\note{le début de la page (jusqu'à « $= z\overset{e}{\otimes}x$ ») est barré de trois traits obliques ; l'argument commence sur une page qui ne figure pas dans ce lot.}
\struck{utilisant la formule}
\[
  \text{\struck{$(a\overset{e}{\otimes}b)\circ(c\overset{e}{\otimes}d) \simeq c\overset{e}{\otimes}b\,(d\wedge a/e)$}} ,
\]
\struck{on trouve}
\[
  \text{\struck{$u\circ w = \bigl(y\overset{e}{\otimes}(-x\wedge y/e)\bigr)\, z\overset{e}{\otimes}z$}}
\]
\[
  \text{\struck{$= -z\otimes x\,(z\wedge y/e) = z\overset{e}{\otimes}x$}}
\]
\note{sous les deux facteurs soulignés $(-x\wedge y/e)$ et $(z\wedge y/e)$ est inscrit « $-1$ ».}

Il existe donc un unique $\rho\in GL(E)$ tel que
\[
  \rho x = y, \quad \rho y = z, \quad \rho z = x
\]
et ce $\rho$ est d'ordre $3$. Il ne change pas \uncertain{quand} on remplace $\{x,y,z\}$ par $(-x,-y,-z)$, il ne dépend que du triplet $\{\bar x,\bar y,\bar z\}$ et d'une orientation de ce triplet. Changer l'orientation, revient à changer $\rho$ en $\rho^{-1}$.

On a aussi
\[
  \det\rho = 1, \quad \operatorname{Tr}\rho = 1, \quad\text{donc}\quad \rho^{-1} = 1-\rho = \rho' \quad (\text{conj. de } \rho)
\]
Considérons
\[
  \tau = x\overset{e}{\otimes}x + y\overset{e}{\otimes}y + z\overset{e}{\otimes}z
\]
On a $\operatorname{Tr}\tau = 0$, $\det\tau = \struck{4}\,3$ (par \ill{} d'\ill{} \ill{} des dét.…) \uncertain{C'est} \ill{} \ill{} qui commute à $\rho$, donc de la forme
\[
  \tau = a\rho + b
\]
d'où, en prenant les traces
\[
  \operatorname{Tr}\tau = a + 2b \quad\text{donc}\quad a = -2b, \text{ donc}
\]
\[
  \tau = \struck{a\rho+2b}\; b(-2\rho+1)
\]
\[
  \det\tau = b^2\det(1-2\rho) \qquad\text{où } \det(1-2\rho) = 3
\]
\note{la lecture « $\operatorname{Tr}\rho = 1$ » est celle du manuscrit ; avec $\det\rho = 1$ et $\rho^3 = 1$, $\rho\neq 1$, on attendrait $\operatorname{Tr}\rho = -1$. La suite ($\rho^{-1} = 1-\rho$, $\det(1-2\rho) = 3$) est cohérente avec $\operatorname{Tr}\rho = 1$, i.e. avec $\rho$ d'ordre $6$ ; le manuscrit dit « d'ordre $3$ ».}

\page{70}
Donc $b^2 = 1$ i.e. $b = \pm 1$, donc
\[
  \tau = \pm(1-2\rho)
\]
\struck{Supposons qu'on ait}
\[
  \tau = 1-2\rho
\]
alors \struck{on ne peut} \ill{} on \ill{} \struck{\uncertain{temps}} $-(1-2\rho)$ (car $2 \cdot 1_k \struck{{}\neq 0}$), donc cela fixe $\rho$ dans la \struck{\uncertain{paire}} $\{\rho, \rho' = 1-\rho\}$. Mais il faut voir lequel\,! \struck{Donc quelle \ill{} \ill{} \uncertain{satisfait} la \ill{},} \struck{\ill{}}
\[
  \text{\struck{$u+v+w = 1-2\rho$}}
\]
Revenant : $M$, on peut dire que pour un triangle $\{u,v,w\}$ de $C^*(e)$, il existe
\[
  \rho\in M^* \text{ unique}
  \left\{
  \begin{array}{l}
    \det\rho = 1 \\
    \rho^3 = 1
  \end{array}
  \right.
  \qquad
  \begin{array}{l}
    \rho \text{ induit une permutation} \\
    \text{circulaire entre les } u, v, w
  \end{array}
\]
et \ill{} \ill{} $\operatorname{Tr}\rho = \operatorname{Tr}\rho' = 1$ \add{\ill{}}
… que $\rho$, ou plutôt \uncertain{si} $\rho'$ pris. Il y a cependant un seul tel $\rho$, tel que
\[
  \boxed{\ u+v+w = 1-2\rho\ }
\]
vérifie \uncertain{ici} cette \uncertain{relation}\,! D'où la structure \struck{\ill{}} orientée : \add{quand} $\{u,v,w\}$ forment un triangle, \struck{\ill{}} le $\rho$ qui \uncertain{vient} de dire, \uncertain{induit} sur $\{u,v,w\}$ une permutation définit une \uncertain{orientation} circulaire canonique.

\marginal{C'est \uncertain{alors} \ill{} $\operatorname{Tr}\rho = 1$ \ill{} \ill{} \ill{} \ill{} \ill{} \ill{} \ill{} $\det\rho = {}$\ill{} \ill{} $\rho^3 = 1$ \ill{} \ill{} \ill{} \ill{}}
\note{note marginale oblique dans la marge gauche ; lecture très incomplète.}

\page{71}
et on a un \struck{bijection} diagramme commutatif

\begin{tikzcd}[column sep=large]
{[E^*/k^* \simeq \mathbb{P}(E)(k)]} \arrow[r, "\sim"] & C^*/k^* \\
E^* \arrow[u] \arrow[r, "\varphi^e"] & C^*(e)\subset C^* \arrow[u]
\end{tikzcd}

\note{au-dessus de la flèche horizontale supérieure : « induit \struck{\ill{}} bijection de schémas » ; au-dessous : « bien connue ». Sous le crochet de gauche : « si $\mathrm{Pic}(k) = 0$ \ill{} », avec une petite flèche. Une flèche en pointillé monte en oblique vers $C^*/k^*$ depuis la droite ; elle n'est pas reproduite.}

On a $\varphi^e(x) = \varphi^e(y)$ ssi $y = \lambda x$ avec $\lambda\in\mu_2(k)$ (i.e. $\lambda\in k^*$, $\lambda^2 = 1$).

Supposons maintenant
\[
  \{\pm 1\} \xrightarrow[\sim]{\text{iso}} \mu_2(k) .
\]
[Si $2$ inv. cela signifie que $\operatorname{Spec} k$ est \uncertain{connexe}. Si $\operatorname{Spec} k$ \uncertain{non} \uncertain{réduit}, $\mathbb{Q}$-connexe. \struck{\ill{}} Si \struck{\ill{}} $2$ \ill{} \uncertain{nul} \ill{}, cela signifie que \ill{} \ill{} \uncertain{composantes} connexes. \struck{\uncertain{Avoir}} \ill{} \ill{} \uncertain{conn.} où $2$ est \uncertain{nul} est $-1$ — i.e. $k = k'\times k''$ avec $2$ \uncertain{inv.} dans $k'$, \struck{\ill{}} \uncertain{produit}, avec \ill{} \ill{} nul dans $k''$, et $\operatorname{Spec}(k'')$ est connexe.]

Alors l'homomorphisme induit
\[
  \psi^e : E^*/\pm 1 \xrightarrow{\ \sim\ } C^*(e)
\]
est bijectif. Cela nous permet de transporter sur $C^*(e)$ la structure

\marginal{\ill{} \ill{} \ill{} dans \ill{} \ill{} \ill{} \ill{} \ill{} \uncertain{comp.} \uncertain{conn.}}
\note{note marginale oblique en regard du crochet ; lecture très incomplète.}

\page{72}
combinatoire $I_e$ sur $E^*/\pm 1$ : structure de \struck{graphe} \add{\uncertain{bi}-polyèdre} \emph{orienté}. Je rappelle que cette structure consiste en

a) Structure de graphe : $\{\bar x,\bar y\}$ forment une arête \struck{\ill{}} ssi $\widetilde{x\wedge y} = e$ ou $-e$ ;

$\longmapsto$ structure triangulaire $\{x,y,z\}$ est un triangle ssi $\{\bar x,\bar y\}$, $\{\bar y,\bar z\}$, $\{\bar z,\bar x\}$ sont des arêtes.

\emph{NB} a) Chaque arête est contenue dans exactement $2$ triangles \add{b) les triangles sont les triplets de la forme $\{\bar x,\bar y,\bar z\}$, $x+y+z = 0$,} \uncertain{est} \emph{orienté}.

b) Chaque triangle \struck{\uncertain{définit} $\{\bar x,\bar y\}$ est \uncertain{divisé} (\ill{}} $\{\bar x,\bar y\}$ des triangles définis l'orientation \struck{de $x\wedge y$} est $x+y+z = 0$, par $x,y,z$ avec $x\wedge y = e$, d'où $y\wedge z = e$, $z\wedge x = e$ \uncertain{distincts}, ssi \add{les deux \uncertain{orientations} d'une arête \uncertain{induites} par celles des deux triangles \uncertain{adjacents} sont \uncertain{opposées}}

c) L'une des arêtes, qui est une \uncertain{incidente} : une arête, \ill{} \ill{} \ill{} combinatoire \ill{} $\mathbb{Z}$ opère \ill{} \ill{} sur laquelle \ill{} donnant structure \ill{} (\ill{} par les \ill{}) \uncertain{induisant} la \ill{} de \ill{} \ill{} \ill{} \ill{} $K$ dessus, d'une opération de $\mathbb{Z}$ \ill{} \add{(\uncertain{inutile})} prolongeant celle de \ill{} arêtes (et \ill{} lesquelles \ill{} des \ill{} \struck{\uncertain{triangles}} ($\longmapsto$ arêtes) sont des $\mathbb{Z}$-torseurs.
\note{les alinéas b) et c) sont très raturés et chargés d'ajouts interlinéaires ; lecture très incomplète. La phrase se poursuit p.~73.}

\page{73}
\emph{NB} Si $s = \bar x$, \struck{\ill{} \ill{}} on définit $T_\lambda = T^s_\lambda$ sur l'ens. des \ill{} \ill{} à $s$, i.e. des $\bar y$ avec $x\wedge y = e$, par
\[
  T_\lambda(\bar y) = \overline{y+\lambda x}
\]
\struck{(\ill{} \ill{} $T_\lambda(y) = ($} (si on remplace $x$ par $-x$, on trouve la \uncertain{même} opération $T_\lambda$). En fait, ici \ill{} l'opér. $T_\lambda$ est induite par la transvection de $E$, $\mathrm{id}+\varphi^e(x\otimes x)$ \struck{$E\xrightarrow{u}$}.

Cette structure se transporte sur $C^*(e)$. Il nous \uncertain{faut} la \uncertain{pouvoir} \ill{} \ill{} l'expliciter : \ill{}

Notons que
\[
  \underbrace{x\overset{e}{\otimes}x}_{=\varphi^e(x)} \circ \underbrace{y\overset{e}{\otimes}y}_{=\varphi^e(y)}
  = y\overset{e}{\otimes}x\,\bigl((x\wedge y)/e\bigr)
\]
\marginal{\emph{NB} $\operatorname{Tr} y\overset{e}{\otimes}x = (x\wedge y)/e$}
donc
\[
  \operatorname{Tr}(u\circ v) = -\bigl((x\wedge y)/e\bigr)^2
\]
\[
  x\wedge y = \pm e \iff \boxed{\ \operatorname{Tr} uv = -1\ }
\]
\note{sous $u$ dans $\operatorname{Tr}(u\circ v)$ est rappelé, en partie biffé, $u = x\overset{e}{\otimes}x$.}
\ill{} les triangles $\{u,v,w\}$ de $C^*(e)$ sont les triples $\{u,v,w\}$ tels que
\[
  \operatorname{Tr} uv = \operatorname{Tr} vw = \operatorname{Tr} wu = -1
\]
Si \struck{\ill{}}
\[
  x+y+z = 0 \qquad x\wedge y = y\wedge z = z\wedge x = e
\]
\[
  u = x\overset{e}{\otimes}x, \quad v = y\overset{e}{\otimes}y, \quad w = z\overset{e}{\otimes}z
\]
\note{la page s'achève sur ces formules ; l'argument s'interrompt ici dans ce lot : le feuillet suivant ouvre une autre liasse, datée d'avril 1986.}

\marginal{On a donc \uncertain{trouvé} sur $C^*$ une structure \ill{} qui \ill{} \ill{} « \ill{} » \ill{} \ill{} \ill{} « \ill{} » \ill{} \ill{} \uncertain{caractéristique} \ill{} ?? \ill{} $\operatorname{Tr} uv = +1$}
\note{note marginale oblique dans la marge gauche ; lecture très incomplète.}

\section{Azumaya rang 4 — avril 1986}
\note{titre et date de la main de l'auteur, seuls inscrits en haut à droite d'un feuillet par ailleurs blanc (page 74 du fonds), qui sert de couverture à la liasse qui suit.}

\page{75}
\note{le feuillet est écrit au verso d'une page dactylographiée qui transparaît ; seule l'écriture manuscrite est transcrite.}
\[
  \delta(xyzt) = \delta(xy(zt)) = Q\bigl(\delta(x),\delta(y),\delta(zt),\delta(xy),\delta(yzt),\delta(ztx)\bigr)
\]
\[
  \text{\struck{$= \delta(x(yz),t)$}}
\]
\note{sous $\delta(yzt)$ et $\delta(ztx)$, soulignés : « remplaçons les pol. en $\delta(y),\delta(z),\delta(t)$ » ; lecture douteuse.}
\[
  = \delta(x(yz)t) = Q\bigl(\delta(x),\delta(yz),\delta(t),\delta(xyz),\delta(yzt),\delta(tx)\bigr)
\]
\[
  = \delta((xy)zt) = Q\bigl(\delta(xy),\delta(z),\delta(t),\delta(xyz),\delta(zt),\delta(txy)\bigr)
\]
\drawing{petit schéma reliant par des traits les mots de trois lettres $yzt$, $ztx$, $xyz$, $txy$ apparus dans les trois expressions.}

On trouve trois expressions polynômiales différentes d'une \ill{} quantité, d'où $2$ relations (\uncertain{quasi}\,! \uncertain{\ill{}}) entre les dix quantités
\[
  \delta(x),\delta(y),\delta(z),\delta(t),\ \delta(xy)\,\delta(xz)\,\delta(xt)\,\delta(yz)\,\delta(yt)\,\delta(zt)
\]
\struck{\ill{} \ill{} \ill{} \ill{}}

$\tau$ \uncertain{normalise} $H\subset G$, $H$ \uncertain{commutatif}, \ill{} que si $h\in H$,
\[
  \boxed{\ \delta(\tau h) \overset{?}{=} \delta(\tau)\ }
\]
\ill{} $\tau^n(h)$ \add{$\overset{\text{déf}}{=} \tau^n h \tau^{-n}$} \uncertain{commute} à $h$ ($n\in\mathbb{Z}$), \ill{} \ill{} \ill{} la formule, en \uncertain{écrivant}. On aimerait l'écrire comme une identité \uncertain{inconditionnelle} \add{(\uncertain{pour} $t\in\mathbb{Z}$)} \struck{\ill{} \ill{}}
\[
  \delta(\tau h) - \delta(\tau) = \text{expression polynomiale en}
\]
\[
  \underbrace{\delta(\{\tau^{n}(h), h\})}_{\text{commutateurs}} \text{ et } \delta(\tau)\,\delta(\tau h),\ \delta(h)
\]
\note{sous l'accolade : « dont tous les termes contiennent un $\delta(\{\ldots\})$ » ; lecture douteuse. À droite, les commutateurs $\{\tau h\tau^{-1}, h\}$ et $\{\tau^{-1}h\tau, h\}$ sont développés : $\tau h\tau^{-1}h\tau h^{-1}\tau^{-1}h^{-1}$ et $\tau^{-1}h\tau h\tau^{-1}h^{-1}\tau h^{-1}$.}
(peut-être suffit-il de $\{\tau h\tau^{-1}, h\}$ et $\{\tau^{-1}h\tau, h\}$)

A-t-on \ill{} \ill{}
\[
  \boxed{\ \delta(\tau.\underbrace{\tau(h)h^{-1}}_{\{\tau,h\}}) \overset{?}{=} \delta(\tau)\ }
  \qquad (\text{si } \tau \text{ normalise } H \text{ commut., } h\in H)
\]

\marginal{C'est idiot\,!, \ill{} \ill{} \ill{} \ill{} \uncertain{expression} \ill{} \ill{} \ill{} \ill{} \ill{} \ill{} \ill{} \ill{} \ill{} \ill{} \ill{} \ill{} \ill{} \ill{} \ill{} \ill{} \ill{} …}
\note{note marginale oblique dans la marge gauche, soulignée « idiot » ; lecture très incomplète.}

\page{77}
\marginal{Description \add{\uncertain{élémentaire}} de $A \simeq M_2(k)$ ($k$ \uncertain{corps}) en termes de $G = A^*/k^*$ (\uncertain{tentatives})}
\note{encadré en haut à droite de la page, en guise de titre. Dans la définition de $\delta(\bar x)$ comme dans la troisième relation ci-dessous, un premier membre illisible est biffé avant la formule retenue.}
\[
  \left[\ x*y = \bar x\,\bar y\,\bar x^{-1}\,\bar y^{-1} \qquad
  \delta(\bar x) = \frac{b^2-4c}{c} \ \right]
  \qquad c = \det x,\ b = \operatorname{Tr} x
\]
\[
  1*1 = 1
\]
\[
  (x*y)(y*x) = 1
\]
\[
  xs*y = x*y \quad\text{si } s, y \text{ commutent}
\]
\[
  [\, x*yt = x*y \quad\text{si } t, x \text{ ——}\,]
\]
\[
  [\, x*y = 1 \quad\text{si } x, y \text{ commutent}\,]
\]
\[
  (x*y)^2 - P(\delta x,\delta y,\delta(xy))\, x*y + 1 = 0
  \qquad P\in\mathbb{Z}[X,Y,Z] \text{ bien déterminé}
\]
\[
  [\,\delta(z), x*y\,] = [\,\delta(z),\delta(z')\,] = 0
\]
\[
  \left\{
  \begin{array}{l}
    \delta(xy) = \delta(yx) \qquad\text{i.e. } \delta(g(x)) = \delta(x) \qquad (x,y,g\in G) \\
    \delta(x) = \delta(x^{-1}) \\
    \delta(xy) + \underbrace{\delta(xy^{-1})}_{=\delta(x^{-1}y)} = (\delta(x)+2)(\delta(y)+2) - 4
  \end{array}
  \right.
\]
\[
  = \delta(x)\delta(y) + 2(\delta(x)+\delta(y))
\]
\note{la troisième relation est cernée ; un trait la relie à « si $x, y$ commutent seulement ??? », et au-dessous : « À vérifier\,! ».}
\[
  \Bigl[\ \boxed{\delta(1) = 0} \Bigr]
\]
$\delta(x^n) = P_n(\delta(x))$, $\delta(x^n y^m) = P_{n,m}(\delta(x),\delta(y),\delta(xy))$ etc. \add{non commutatifs…} \ill{} d'autres monômes non commutatifs ont une valeur pol.

\emph{Pour} $x_1, x_2, x_3, x_4\in G$, il doit y avoir une relation à coefficients dans $\mathbb{Z}$ entre les
\[
  \delta(x_i)\ (1\leq i\leq 4), \qquad \delta(x_i x_j)\ (1\leq i<j\leq 4)
\]
(dix quantités, \struck{\ill{} \ill{} \ill{}} « \uncertain{indépendantes} » de \uncertain{celle} \uncertain{précédente} …]
\note{tout ce paragraphe est marqué d'un double trait vertical dans la marge.}
\[
  \boxed{\ \delta(xyz) \overset{?}{=} Q\bigl(\delta(x),\delta(y),\delta(z),\delta(xy),\delta(yz),\delta(zx)\bigr)\ }
\]
(pol. univ. à coefficients \uncertain{entiers} en six variables, déterminé)
\struck{Ceci donne $\delta(x^n)$ comme un pol.}
\[
  \delta(x^2) = (\delta(x)+2)^2 - 4 = \delta(x)^2 + 4\delta(x)
  \qquad [\text{cas particulier de } \delta(xy)+\delta(xy^{-1}) = \ldots]
\]
\[
  \delta(x^n) = P_n(\delta(x)) \qquad n\in\mathbb{Z} \qquad P_n = P_{-n}
\]
\[
  \delta(x^n y^m) = P_{n,m}(\delta(x),\delta(y),\delta(xy))
\]
\[
  \delta(\underbrace{M(x_1,\ldots,x_n)}_{\substack{\text{monôme non commutatif}\\ \text{à exposants entiers}\,\in\mathbb{Z}}})
  = P_M\bigl((\delta(x_i)),(\delta(x_i x_j))\bigr)
\]
\note{sous $P_M$ : « pol. à coeff. dans $\mathbb{Z}$ (pas uniquement déterminé si $n\geq 4$) ». Avant $M(x_1,\ldots,x_n)$ le mot « monôme » est biffé puis réécrit dessous.}

\page{80}
\emph{Azumaya rang 4} (I)

Avril 86

1. Soient $k$ anneau, $A$ algèbre d'Azumaya sur $k$ de rang $4$ (donc loc. iso. à $M_2(k)$). On a alors
\[
  (1)\quad
  \left\{
  \begin{array}{lll}
    \operatorname{Tr} : A\longrightarrow k & k\text{-linéaire} & \\
    \det : A\longrightarrow k & \text{multiplicatif} & \text{, quadratique}
  \end{array}
  \right.
\]
tels que $\forall u\in A$, on ait
\[
  (2)\quad u^2 - \operatorname{Tr}(u).u + \det(u) = 1
\]
(Hamilton-Cayley). De plus, on a
\[
  (3)\quad \det(u+v) = \det u + \det v + (\operatorname{Tr}u\operatorname{Tr}v - \operatorname{Tr}(uv))
\]
\note{dans (2), le second membre se lit « $1$ » ; on attendrait $0$.}
(i.e.
\[
  (u,v)\longmapsto \Phi(u,v) \overset{\text{déf}}{=} \operatorname{Tr}u\operatorname{Tr}v - \operatorname{Tr}(uv)
\]
est la forme bil. ass. à la forme quad. $\det$).

\emph{Remarques} On peut s'intéresser aux $k$-alg. quelconques $A$ pour lesquelles on a (1) (2) (3). Si $A$ est proj. de t.f. disons et $2$ inv. dans $k$, alors (1) (2) $\Rightarrow$ (3), car il suffit de le voir pour $u = v$ (deux formes bil. \emph{sym.} qui coïncident pour $u = v$ coïncident)\struck{, et alors cela se \uncertain{résulte de} \uncertain{réduit} à une vérification dans les $k[u] = k\oplus k.u$ quadratiques sur $k$…]}. La relation pertinente est
\[
  (4)\quad (\operatorname{Tr}u)^2 - \operatorname{Tr}(u^2) = 2\det u ,
\]
et elle résulte de (2) en l'écrivant sous la forme
\[
  (2')\quad u^2 = \operatorname{Tr}(u)u - \det(u)
\]
et portant cette valeur de $u^2$ dans (4), compte tenu…
\note{la phrase se poursuit au-delà de la page 80, hors de ce lot.}

\end{document}
